\documentclass[preprint,12pt,nonatbib]{elsarticle}

\journal{Consciousness and Cognition}

\usepackage[T1]{fontenc}
\usepackage[utf8]{inputenc}
\usepackage{lmodern}
\usepackage{microtype}
\usepackage{booktabs}
\usepackage{amsmath,amssymb}
\usepackage{hyperref}
\usepackage{xurl}
\usepackage[style=apa,backend=biber,natbib=true,sortcites=true]{biblatex}
\DeclareLanguageMapping{english}{english-apa}
\addbibresource{references.bib}

\hypersetup{colorlinks=true,linkcolor=blue,citecolor=blue,urlcolor=blue}
\setlength{\parindent}{0pt}
\setlength{\parskip}{6pt plus 2pt minus 1pt}

\begin{document}

\begin{center}
{\Large\bfseries Fitted-Operator Geometry as a Complementary Descriptive Axis for Brain-State Discrimination in Human iEEG and Scalp EEG\par}
\vspace{6pt}
{\normalsize Phillip Peterkin\textsuperscript{1}\par}
\vspace{3pt}
{\footnotesize \textsuperscript{1}Independent Researcher, Albany, Oregon, United States\par}
\vspace{3pt}
{\footnotesize ORCID: 0009-0006-4525-6685\par}
\end{center}

\section*{Abstract}
Criticality, complexity, and spectral-power measures compress neural recordings into useful scalar summaries, but they do not directly describe the geometry of fitted multivariate operators. We evaluated whether fitted first-order vector autoregressive (VAR(1)) operators provide complementary descriptive coordinates across human intracranial electroencephalography (iEEG), propofol electroencephalography (EEG), and sleep EEG datasets. In the primary iEEG band comparison, a threshold-derived branching statistic was lower in high-gamma than broadband activity ($\sigma_{HG}=0.9735$ versus $\sigma_{BB}=0.9908$; paired $t(17)=-5.74$, $p=8.9\times10^{-6}$). A historical geometry-proximity statistic, defined from eigenvector overlap and minimum eigenvalue spacing, was associated across subjects with the branching statistic ($r=0.860$, $p=4.8\times10^{-6}$); this historical statistic is explicitly distinct from the current Near-Degeneracy (ND) construction. In scalp EEG, minimum eigenvalue spacing distinguished propofol sedation from wakefulness ($d=0.71$) and rapid-eye-movement (REM) from N3 sleep ($d=-2.51$). A historical four-feature state-classification artifact achieved leave-one-subject-out area under the receiver operating characteristic curve (AUC) $=0.9125$ for awake versus propofol, while sleep contrasts reached high within-cohort separation; those historical feature bundles contain the legacy proximity statistic and are not presented as validation of the current ND score. Before N2-to-N3 transitions, spectral radius changed across the 120-second pre-boundary window ($p=0.0014$; non-overlapping validation $p=0.032$). Shared-subspace, alternative-spacing, regularization, non-overlap, and surrogate analyses constrain interpretation. All findings are conditional on the declared estimator and preprocessing pipeline and are treated as properties of fitted local linearizations rather than identified neural mechanisms.

\textit{Keywords:} vector autoregression; operator geometry; eigenvalue spacing; eigenvector conditioning; criticality; anesthesia; sleep; reproducible computational neuroscience

\section{Introduction}
Neural-state analyses commonly summarize activity through oscillatory power, avalanche statistics, complexity, and temporal correlations \citep{Beggs2003,Hardstone2012,Lombardi2023}. These quantities are useful but leave open a complementary question: what structure is visible in a fitted multivariate dynamical operator under a fixed estimation pipeline?

We address that question using sliding-window VAR(1) models. The fitted coefficient matrix is not interpreted as the true neural generator. Instead, eigenvalue and eigenvector summaries are treated as descriptive properties of a local linear approximation. This distinction is important because preprocessing, dimensionality reduction, finite-window estimation, and autocorrelation can all affect the resulting operator geometry.

The public repository accompanying this report is designed to make those dependencies inspectable. It separates raw-data analysis, machine-readable result artifacts, mathematical unit tests, manuscript-to-result audit tests, robustness analyses, and historical quantities whose interpretation changed during review. In particular, an early proximity statistic stored as \texttt{ep\_score} motivated later work but is not mathematically identical to the current ND score. The two are therefore reported separately throughout this version.

The primary questions are: (1) whether high-gamma and broadband signals differ under the declared criticality-related estimator; (2) whether fitted-operator summaries show reproducible state structure in anesthesia and sleep; (3) whether geometry carries descriptive information not reducible to simple power shifts; and (4) which conclusions survive robustness and falsification tests.

\section{Results}

\subsection{High-gamma and broadband differ under the declared branching-related estimator}
Within the COGITATE iEEG cohort, high-gamma activity produced a lower branching statistic than broadband activity under the same analysis pipeline ($\sigma_{HG}=0.9735$ versus $\sigma_{BB}=0.9908$; paired $t(17)=-5.74$, $p=8.9\times10^{-6}$). A mixed-effects band contrast was also significant (coefficient approximately $-0.0173$, $p=9.3\times10^{-9}$). Lempel--Ziv complexity moved in the opposite direction, with higher values in high-gamma, while detrended fluctuation analysis was higher in broadband. The avalanche exponent difference did not reach the conventional group threshold. These dissociations argue against treating the four criticality-adjacent measures as a single interchangeable axis.

The branching quantity is described here as a criticality-related statistic under a thresholded-event estimator. It is not treated as an estimator-independent measurement of the latent neuronal branching parameter.

\subsection{Historical cross-subject geometry associations are retained with their original statistic}
The repository contains a historical per-window proximity quantity
\begin{equation}
P_t = \frac{o_t}{g_t + 10^{-10}},
\end{equation}
where $o_t$ is the absolute overlap of the closest eigenvector pair and $g_t$ is the minimum eigenvalue spacing. Across 18 iEEG subjects, the subject-level historical proximity summary was associated with the branching statistic ($r=0.860$, $p=4.8\times10^{-6}$), and leave-one-subject-out jackknife values ranged from approximately 0.792 to 0.889. Minimum spacing itself was negatively associated with the branching statistic ($r=-0.588$, $p=0.0102$). Lempel--Ziv complexity was negatively associated with the historical proximity summary ($r=-0.684$, $p=0.00174$) and positively associated with minimum spacing ($r=0.745$, $p=3.9\times10^{-4}$).

These values are preserved because they are part of the computational history. They are \emph{not} relabeled as empirical results of the current ND construction. The distinction is enforced by repository documentation and regression tests.

\subsection{Propofol and sleep show distinct fitted-operator geometry}
In the propofol EEG cohort, spectral radius increased slightly under sedation (approximately 0.9980 awake versus 1.0025 sedated; $d=-1.66$, with the sign reflecting the awake-minus-sedated convention). Minimum eigenvalue spacing decreased from approximately 0.0122 to 0.0114 ($p=0.0052$, $d=0.71$). Eigenvector condition number also increased under sedation, whereas effective rank did not show a strong primary effect. The spacing effect persisted under shared-subspace estimation and alternative nearest-neighbor spacing summaries.

Sleep showed a different organization. Awake and N3 did not differ reliably in the headline minimum-spacing statistic, whereas REM displayed tighter spacing. The N3-versus-REM contrast was large ($d=-2.51$, $p\approx2.4\times10^{-5}$), and awake also differed from REM. Because the sleep sample contains only ten subjects, the magnitude and perfect separation observed in some within-cohort classification contrasts are treated as small-sample results rather than population-level performance guarantees.

\subsection{Historical state-classification artifact}
The checked-in \texttt{geometry\_brain\_states.json} artifact evaluates a four-feature bundle containing minimum spacing, eigenvector condition number, spectral radius, and the historical \texttt{mean\_ep\_score} proximity quantity. Earlier code labeled that fourth column \texttt{nd\_score}. That label was semantically incorrect and is retained only inside historical artifacts for provenance.

Using leave-one-subject-out classification with fold-local standardization, the historical feature bundle achieved AUC $=0.9125$ for awake versus propofol in the 20-subject cohort. The same historical pipeline yielded perfect AUC in the N3-versus-REM sleep contrast for the ten subjects available. These values demonstrate that the historical fitted-geometry feature bundle contains state-discriminative information within the declared cohorts. They do not establish external generalization, and they do not validate the current ND definition.

\subsection{Temporal precedence is interpreted as a pre-boundary association}
Before scored N2-to-N3 transitions, spectral radius changed systematically across the preceding 120 seconds (group $p=0.0014$; nine of ten subjects in the reported direction). The effect survived a non-overlapping-window validation ($p=0.032$). Minimum spacing showed a supportive pre-transition trend but did not survive the same stricter non-overlap criterion and is not promoted to an equally robust result.

Because sleep stages are scored in discrete epochs and the analysis is observational, temporal precedence is interpreted as a pre-boundary biomarker under this pipeline, not evidence that the fitted geometry causally controls the transition.

\subsection{Robustness and falsification constrain the interpretation}
Several analyses were designed specifically to test alternative explanations. Shared-subspace principal component analysis (PCA) reduced concern that state-specific coordinate systems create the headline scalp effects. Alternative spacing summaries reduced dependence on an extreme-value minimum-gap statistic. A ridge-regularization sweep showed that the propofol spectral-radius contrast was not confined to one narrow least-squares conditioning regime. Non-overlapping windows preserved the primary propofol effects at the subject level.

Phase-randomized surrogate analyses produced an important negative result: absolute spectral-sensitivity magnitudes in real propofol data did not exceed the surrogate group mean in the tested subset ($p=0.23$). Consequently, absolute sensitivity magnitude is not interpreted as a neural-specific marker. This negative result is retained in the public record rather than removed from the narrative.

\section{Methods}

\subsection{Datasets}
Four external datasets are represented in the public analysis record: COGITATE Experiment 1 iEEG for the primary intracranial analyses; OpenNeuro ds005620 for propofol EEG; ANPHY-Sleep for sleep-state and pre-transition analyses; and OpenNeuro ds004752 for a secondary stereoelectroencephalography consistency/generalization analysis. Raw recordings are not redistributed in this repository. Dataset roles, subject counts, exclusions, and acquisition links are recorded in \texttt{REPLICATION\_AND\_DATA\_PROVENANCE.md} and \texttt{data/README\_data.md}.

\subsection{Preprocessing and fitted operators}
The intracranial analysis compares high-gamma (70--150 Hz) and broadband (1--200 Hz) signal definitions. Scalp analyses use a lower-frequency broadband preprocessing path appropriate to the source recordings, followed by current-source-density transformation and PCA under the declared configuration. Shared-subspace PCA is used as a robustness control for cross-state comparisons.

Within each analysis stream, a VAR(1) model is fitted in 500 ms windows with 100 ms step size. For a multivariate signal $\mathbf{x}_t$, the local model is
\begin{equation}
\mathbf{x}_{t+1} = \mathbf{A}\mathbf{x}_t + \boldsymbol{\varepsilon}_t.
\end{equation}
The coefficient matrix $\mathbf{A}$ is a fitted descriptive operator, not an identified biological transition matrix.

For each window, the primary primitive geometry quantities include spectral radius
\begin{equation}
\rho(\mathbf{A}) = \max_i |\lambda_i|,
\end{equation}
minimum eigenvalue spacing
\begin{equation}
g = \min_{i<j}|\lambda_i-\lambda_j|,
\end{equation}
and eigenvector condition number
\begin{equation}
\kappa(\mathbf{V}) = \|\mathbf{V}\|_2\,\|\mathbf{V}^{-1}\|_2.
\end{equation}
Minimum spacing is dimension-dependent and is therefore interpreted comparatively under fixed dimensional and preprocessing settings.

\subsection{Historical proximity statistic and current Near-Degeneracy score}
Two different composite quantities exist in the project history and must not be conflated.

The historical proximity statistic is
\begin{equation}
P_t = \frac{o_t}{g_t + 10^{-10}},
\end{equation}
where $o_t$ is closest-pair eigenvector overlap. Historical JSON fields such as \texttt{ep\_score} and \texttt{mean\_ep\_score} refer to this quantity.

The current Near-Degeneracy score begins with paired valid windows and defines
\begin{align}
c_t &= -\log_{10}(g_t+\epsilon),\\
k_t &= \log_{10}(\kappa(\mathbf{V}_t)+\epsilon),
\end{align}
with $\epsilon=10^{-12}$. The two feature series are standardized within the declared analysis unit to produce $\tilde c_t$ and $\tilde k_t$. The score is the projection onto the first principal component of the paired standardized feature matrix,
\begin{equation}
ND_t = w_c\tilde c_t + w_k\tilde k_t,
\end{equation}
where $(w_c,w_k)$ is the unit-norm first-principal-component loading vector. Principal-component sign is arbitrary, so same-sign loadings are oriented positive. If the leading loadings have opposite signs or a zero loading, the public implementation fails loudly rather than silently redefining the construct.

Because the current ND inputs are standardized within the supplied analysis unit, the simple mean of ND within that same unit is approximately zero by construction. A subject-mean cross-subject correlation therefore requires a prospectively specified aggregation/normalization rule and cannot be substituted for the historical $P$ statistic. No such substitution is made in this report.

\subsection{Classification and leakage control}
Historical state-classification analyses use subject-condition summaries rather than treating overlapping windows as independent observations. Leave-one-subject-out splits preserve subject identity. Standardization is fitted on training observations inside each fold and then applied to the held-out subject. The historical state-space feature bundle contains the legacy proximity statistic, even where an older JSON field calls that feature \texttt{nd\_score}. Current schema documentation records this mismatch explicitly.

\subsection{Statistical inference}
Primary analyses use subject-level paired tests, effect sizes, correlations, and mixed-effects models as appropriate. False-discovery-rate correction is applied within declared test families. Overlapping windows are not treated as independent subjects. Circular-shift or surrogate procedures are used where preservation of temporal structure is required. The public result artifacts retain both significant and null results relevant to interpretation.

\section{Software-engineering and reproducibility controls}
The repository is an installable Python package (\texttt{operator-geometry}, version 1.0.0) with a canonical configuration at \texttt{code/config.yaml}. Continuous Integration installs the package from a clean runner, checks dependency consistency, runs Ruff static checks, builds a wheel, and executes the portable test suite. The test suite includes mathematical known-answer tests, synthetic fitted-system tests, artifact-to-manuscript checks, public scientific-contract checks, and analysis-module tests.

The public architecture separates loaders, preprocessing, model fitting, feature extraction, statistical analysis, provenance, results, and manuscript auditing. Data roots are supplied through documented environment variables rather than hardcoded private paths. Historical scripts and JSON keys are retained where needed for provenance, but the canonical public path is documented separately.

A dedicated \texttt{PUBLIC\_AUDIT.md} records unresolved release gates. A dedicated \texttt{SCIENTIFIC\_INTEGRITY.md} records subject-boundary, leakage, aggregation, exploratory/confirmatory, and result-overwrite rules. The purpose of these controls is not to claim that the pipeline is error-free; it is to make errors detectable, localizable, and correctable.

\section{Limitations}
\begin{enumerate}
\item Fitted VAR(1) operators are descriptive local linearizations. Their geometry does not identify the underlying neural mechanism.
\item The branching statistic depends on event thresholding, filtering, sampling, and estimator choice and is not treated as a direct measurement of a latent branching-process parameter.
\item Minimum eigenvalue spacing is dimension-dependent and sensitive to finite-sample estimation.
\item The historical proximity statistic and the current ND score are distinct. Historical cross-subject correlations and state-classification artifacts cannot be silently promoted into current-ND validation.
\item The propofol and sleep cohorts are modest in size, particularly the ten-subject sleep cohort. Within-cohort classification performance is not an external validation estimate.
\item Some raw-data reproduction gates remain open in the public audit, including end-to-end confirmation of the canonical broadband path against the historical checked-in result.
\item Surrogate analyses show that some sensitivity quantities can arise from autoregressive fitting to autocorrelated signals. Those quantities are interpreted comparatively rather than as neural-specific biomarkers.
\end{enumerate}

\section{Data and code availability}
The iEEG data are available from the COGITATE Consortium. The secondary stereoelectroencephalography dataset is OpenNeuro ds004752. The propofol EEG dataset is OpenNeuro ds005620. The sleep dataset is available through the ANPHY-Sleep OSF record. Code, configuration, machine-readable result artifacts, tests, and audit documentation are available at \url{https://github.com/Phillip-Peterkin/Eigenvalue-and-eigenvector-geometry}.

\section{Author contributions}
P.P. conceived the study, designed the analysis pipeline, performed the analyses, built the public reproducibility and audit structure, and wrote the manuscript.

\section{Competing interests}
The author declares no competing interests.

\clearpage
\printbibliography
\end{document}
